\documentclass[11pt,a4paper]{article}
\usepackage[T1]{fontenc}
\usepackage[utf8]{inputenc}
\usepackage[portuguese]{babel}
\usepackage[margin=2.2cm]{geometry}
\usepackage{enumitem}
\usepackage{titlesec}
\usepackage{xcolor}
\usepackage[hidelinks]{hyperref}
\usepackage{booktabs}
\usepackage{amssymb}

\definecolor{cinza}{gray}{0.35}
\titleformat{\section}{\large\bfseries}{}{0pt}{}[\vspace{-6pt}\rule{\textwidth}{0.4pt}]
\titlespacing{\section}{0pt}{12pt}{5pt}
\newcommand{\lin}[2]{\noindent\textbf{#1}\hfill{\color{cinza}\small #2}\par}
\setlist[itemize]{leftmargin=1.2em,itemsep=1pt,topsep=2pt,parsep=0pt}
\pagestyle{empty}

\begin{document}

% CC BY 4.0 · Copyright (c) 2026 Aarão Melo Lopes · ver LICENSE na raiz.

\begin{center}
{\LARGE\bfseries Aarão Melo Lopes}\\[2pt]
{\color{cinza}\emph{Engenheiro Eletricista \,|\, Matemática Aplicada \& Engenharia de Sistemas}}\\[2pt]
{\small
\href{mailto:aarao.melo.lopes@gmail.com}{aarao.melo.lopes@gmail.com}
\quad·\quad
\href{https://goldenkingdom.patriatechnology.com}{goldenkingdom.patriatechnology.com}}
\end{center}

\vspace{-2pt}
\section*{O que faço}

A dualidade é a memória da divisão: guardar a metade que se ia deitar fora, e tornar a
divisão reversível. Em termos técnicos --- um \textbf{corpo dual} $(V,V^{*},J)$ e duas leis
bem tipadas ($1^{\dagger}=-1$, $T^{\dagger}=-T$). Cada afirmação vem com um teste em C que a
pode desmentir.

\smallskip
Três volumes, um papel cada: a \textbf{Teoria} afirma; o \textbf{Catálogo} é o
\emph{embaixo} --- bestiário onde cada corpo tem régua e dinâmica, e
\emph{o que não é citado não é testado}; o \textbf{Enredo} \emph{não inventa matemática ---
conta-a}, na linguagem do xadrez. Começou na graduação (hipercomplexos), atravessou o
IME-USP, e é hoje 58 + 446 + 360\,pp., 306 testes a zero e seis sistemas.

\smallskip
Dois factos que ligam o percurso:
\begin{itemize}
\item Um contador de bits de 33 linhas, anterior à teoria, foi reencontrado por ela e medido
propriedade a propriedade.
\item A tabela de comutação do TCC de 2018 (inversor de três níveis) coincide com a ISA
mínima do catálogo: uma instrução \texttt{MOVE}(destino, sentido), sem estado auxiliar.
\end{itemize}

\smallskip
\noindent\textbf{O que já se verifica hoje (sem adopção externa ainda):}
repositório público, documentos com páginas fixas, bateria a zero, publicação com portão.
Não há ainda citações nem utilizadores em produção fora da bancada.

\lin{Em que pé está cada coisa}{}

\begin{center}\footnotesize
\begin{tabular}{@{}p{2.5cm}p{7.6cm}p{4.4cm}@{}}
\toprule
\textbf{estado} & \textbf{o que é} & \textbf{onde} \\
\midrule
\textbf{demonstrado} &
21 enunciados com prova escrita &
\emph{Teoria}, 58\,pp. \\[2pt]
\textbf{medido} &
306 medidores em C; a bateria sai a zero &
\emph{Catálogo} e \texttt{tests/} \\[2pt]
\textbf{implementado} &
\texttt{MOVE}; banco sem \emph{malloc}; assistente sem modelo
externo; compilador e motor de xadrez &
\texttt{banco/}, \texttt{broca-so}, \texttt{chess} \\[2pt]
\textbf{proposto} &
liga grafeno--estanho; Headjack; CristalChain; jogo (GDD) &
propostas adiante \\
\bottomrule
\end{tabular}
\end{center}

\section*{Produção}

\lin{Uma teoria da dualidade, em três volumes}{2026, em curso}
\begin{itemize}
\item \textbf{Teoria} (58\,pp.) --- 21 enunciados e 21 provas; uma definição --- o corpo dual.
\item \textbf{Catálogo} (446\,pp.) --- corpos um a um; medidores citados; GDD do reino;
CristalChain, ISA, liga, Headjack.
\item \textbf{Enredo} (360\,pp.) --- mito re-encenado a cada partida: Operador de duas mãos;
Cristal (23 lances) e Quasicristal (128 cantos); vitória = \emph{mate que não mata},
resíduo zero com o relógio a correr. $\sim$30 partes, $\sim$150 capítulos.
\item \textbf{306 medidores em C} --- 306 verdes, 0 falhas; 115 só com aritmética inteira.
\end{itemize}

\smallskip
\lin{\emph{O sistema dos números hipercomplexos} --- livro próprio}{graduação}
Semente: $\hat z^{\,2}=-1$ e $z^{-1}=\bar z/|z|^2$ --- hoje as duas leis.

\smallskip
\lin{\emph{Álgebra Bilinear} --- notas próprias}{graduação}
$f(u,v)=[u]^{T}_{B}A[v]_{B}$ --- a matriz como leitura numa base.

\smallskip
\lin{Seis sistemas construídos}{abr.--ago.\ 2026}
Compilador, VM, núcleo, \texttt{broca-so}, motor de xadrez, corpo teórico.
\\[1pt]
{\small\color{cinza}C · Makefile · Shell · POSIX · Git · \texttt{MOVE} · sem \emph{malloc}
· \texttt{\_start}}

\section*{Formação}

\lin{Mestrado em Ciência da Computação --- IME-USP}{2020 -- 2022}
Orientador: Prof.\ Dr.\ Ronaldo Fumio Hashimoto. Bolsista CAPES.
\textbf{72 créditos} (exigência: 44); conceito máximo em oito das nove.
O corpo teórico e os sistemas substituíram a dissertação formal --- disponível sob consulta.

\smallskip
\noindent{\small Análise de Algoritmos · Teoria dos Grafos · Representações Hierárquicas ·
Processamento de Imagens usando Grafos · Visão e Processamento de Imagens ·
Modelos Probabilísticos em Grafos · Computação Paralela e Distribuída · Computação Gráfica.}

\smallskip
\lin{Extensão --- IME-USP}{verão 2021}
Cálculo no $\mathbb{R}^n$ e Espaços Métricos (60\,h cada); frequência 100\,\%.

\smallskip
\lin{Bacharelado em Engenharia Elétrica --- UFRR}{2011 -- 2018}
Média \textbf{8{,}66} · 4\,201 horas.
Nota 10 em: Cálculo I · Cálculo Diferencial e Integral I · Eletromagnetismo I ·
Física Experimental III · Conversão de Energia II · Fundamentos de Economia.
Acima de 9: Estatística · Química Geral · Cálculo Numérico · Métodos Matemáticos ·
Geometria Analítica · Álgebra Linear · Geração de Energia.

\section*{Trabalho de conclusão e docência}

\lin{\emph{Controle Direto de Torque do Motor de Indução Trifásico}}{UFRR, 2018}
123\,pp. Orientadora: Prof.\textsuperscript{a} Dr.\textsuperscript{a} Susset Guerra Jiménez.
Nota 10 da banca, recomendação de publicação.
\\[1pt]
{\small Cada fase $\{-1,0,+1\}$: $27$ estados, 19 vetores, seis pares redundantes --- a
redundância do ponto médio \emph{é} a dualidade; o DTC sem histórico é o \texttt{MOVE}
da Proposta~I.}

\smallskip
\lin{Monitoria --- 473 horas}{UFRR, 2013 e 2016--17}
Cálculo I (108\,h) e Álgebra Linear I no curso de Matemática (365\,h).

\section*{Proposta I --- engenharia}

\lin{\emph{Headjack} + liga grafeno--estanho}{proposta aberta}

Interface sem linguagem intermédia: formalismo (par leitura/escrita, régua, dinâmica) já
escrito; máquina \texttt{MOVE} já corre.
\begin{itemize}
\item \textbf{Headjack} (\texttt{tests/headjack.c}) --- leitura não invasiva (Faraday, Lenz,
Poynting); limites do operador corrente$\to$campo (núcleo; o que não tem dual não atravessa).
\item \textbf{Liga} (\texttt{tests/liga.c}) --- grafeno + estanho: percolação, salto
$\sim 10^{14}$, janela de casamento; modelo calculado, bancada por fazer.
\end{itemize}

\textbf{O que preciso do parceiro:} bancada/indústria para a liga; software para agentes
em produção (\texttt{broca-so} já orquestra; falta produção).

\section*{Proposta II --- jogo e cinema}

\lin{O Enredo --- a partida sem fim}{disponível para colaboração}

Não é campanha com roteiro: é mito que cada partida re-encena. Semente: o
\emph{gambito recusado} --- quem aceita ganha material e paga o centro; quem recusa fica
largo; a oferta fica de pé. Duas metades (terra/floresta/mar e Fortaleza/Alfândega/
Telescópio/Fábrica). GDD no catálogo: dois zeros no HUD (resíduo e relógio); vitória =
$\Gamma=0$ com o pulso a correr --- não o flatline.

\smallskip
\textbf{O que preciso do parceiro:} equipa de jogo ou cinema. Trago o mundo; aceito
qualquer posição.

\section*{Proposta III --- registo distribuído}

\lin{\emph{CristalChain} --- livro-razão cuja cifra conserva}{proposta aberta}

No enredo é o \emph{ouro externo}: acrescenta sem tirar a ninguém. Na álgebra: cifra
forçada pelo determinante --- ouro branco ($\det=+1$) conserva; metais ($\det=-1$)
anti-conservam. Paragem interna (órbita finita / período de Pisano), não gás externo.

\smallskip
\textbf{Existe:} álgebra e medidores (\texttt{ouro\_branco.c}, \texttt{smartcontract.c}).
\textbf{Não existe:} rede, nós, consenso.
\emph{Não é blockchain; não resolve dupla despesa.}

\smallskip
\textbf{O que preciso do parceiro:} sistemas distribuídos / consenso.

\section*{Como trabalho e como colaborar}

\begin{itemize}
\item Meço em vez de argumentar --- cada afirmação com um teste que a pode derrubar.
\item \textbf{Presencial:} grupo, bancada e prazo --- escreva-me.
\item \textbf{Remota:} CC BY 4.0 / MIT; pode avançar sem combinar.
\end{itemize}

\medskip
\begin{center}\small
\textbf{Documentos e código}\\[4pt]
\begin{tabular}{@{}ll@{}}
teoria &
\href{https://goldenkingdom.patriatechnology.com/docs/teoria.pdf}{\texttt{.../docs/teoria.pdf}}
\quad (58\,pp.) \\
catálogo &
\href{https://goldenkingdom.patriatechnology.com/docs/catalogo.pdf}{\texttt{.../docs/catalogo.pdf}}
\quad (446\,pp.) \\
enredo &
\href{https://goldenkingdom.patriatechnology.com/docs/enredo.pdf}{\texttt{.../docs/enredo.pdf}}
\quad (360\,pp.) \\
código &
\href{https://goldenkingdom.patriatechnology.com/repo.git}{\texttt{git clone \ldots/repo.git}} \\
\end{tabular}
\\[4pt]
{\footnotesize\color{cinza}CC BY 4.0 (textos) · MIT (código)}
\end{center}

\end{document}
